\begin{table}[htbp]
\caption{Rajasthan: Balance Tests for Quota Assignment}
\label{tab:balance_raj_2005_2010}
{\centering\scriptsize
\begin{tabular}{@{}lrrr rrr rrr@{}}
\toprule
& \multicolumn{3}{c}{2005$\rightarrow$2010} & \multicolumn{3}{c}{2010$\rightarrow$2015} & \multicolumn{3}{c}{2015$\rightarrow$2020} \\
\cmidrule(lr){2-4} \cmidrule(lr){5-7} \cmidrule(lr){8-10}
Variable & Quota & Open & p & Quota & Open & p & Quota & Open & p \\
\midrule
Population & 4088.1 & 4118.3 & 0.475 & 4069.8 & 4126.0 & 0.179 & 4112.5 & 4103.0 & 0.822 \\
Female (\%) & 48.4 & 48.3 & 0.213 & 48.4 & 48.4 & 0.540 & 48.4 & 48.4 & 0.246 \\
SC (\%) & 18.3 & 18.1 & 0.529 & 18.5 & 17.7 & 0.014 & 18.0 & 18.0 & 0.876 \\
ST (\%) & 14.0 & 14.4 & 0.556 & 14.3 & 14.7 & 0.466 & 14.6 & 14.6 & 0.952 \\
Medical facility (\%) & 84.0 & 84.9 & 0.438 & 85.2 & 84.0 & 0.258 & 84.3 & 85.3 & 0.367 \\
Maternity homes/1k & 0.0 & 0.0 & 0.586 & 0.0 & 0.0 & 0.729 & 0.0 & 0.0 & 0.815 \\
Maternal/child centers/1k & 0.0 & 0.0 & 0.064 & 0.0 & 0.0 & 0.589 & 0.0 & 0.0 & 0.519 \\
Primary schools/1k & 1.1 & 1.1 & 0.702 & 1.1 & 1.1 & 0.057 & 1.1 & 1.1 & 0.068 \\
Middle schools/1k & 0.4 & 0.4 & 0.038 & 0.4 & 0.4 & 0.658 & 0.4 & 0.4 & 0.676 \\
Banking (\%) & 21.8 & 21.8 & 0.979 & 21.2 & 22.2 & 0.414 & 22.1 & 21.6 & 0.682 \\
Power supply (\%) & 98.6 & 97.9 & 0.063 & 98.3 & 98.3 & 0.993 & 98.2 & 98.4 & 0.609 \\
Mud road (\%) & 85.9 & 86.7 & 0.492 & 86.4 & 86.5 & 0.914 & 85.9 & 87.0 & 0.270 \\
\midrule
N & \multicolumn{3}{c}{4980} & \multicolumn{3}{c}{4605} & \multicolumn{3}{c}{4514} \\
Joint p (LR) & \multicolumn{3}{c}{0.247} & \multicolumn{3}{c}{0.159} & \multicolumn{3}{c}{0.598} \\
Joint p (block) & \multicolumn{3}{c}{0.019} & \multicolumn{3}{c}{0.132} & \multicolumn{3}{c}{0.301} \\
\bottomrule
\end{tabular}
\par}

\vspace{0.5ex}
\parbox{\linewidth}{\scriptsize \emph{Notes:} Balance tests for quota assignment across all Rajasthan panels. ``Quota'' = GPs reserved for women in base year; ``Open'' = GPs not reserved. Covariates from the 2001 Census Village Directory (SHRUG). Availability is the percentage of GPs with an observed facility in at least one linked village; absence requires complete observed coverage. Counts per 1k use GP population. Maternity/child centers are maternity and child welfare centers. Water-source code sums are excluded because they do not identify facility counts or availability. p-values from t-tests. Joint p (LR) is a descriptive pooled logistic likelihood ratio test. Joint p (block) is a Wald test with district-samiti/block fixed effects and standard errors clustered at that level. These are balance diagnostics, not randomization inference.}
\end{table}
